
\documentclass[a4paper,12pt]{article}

% Font e lingua
\usepackage[T1]{fontenc}
\usepackage[utf8]{inputenc}
\usepackage[italian]{babel}
\usepackage{lmodern}
\usepackage{setspace}
\usepackage{amsmath}
\usepackage[protrusion=true,expansion=false]{microtype}

% Impaginazione
\usepackage[
  a4paper,
  left=2.5cm,
  right=2.5cm,
  top=3cm,
  bottom=3cm
]{geometry}

% Interlinea
\usepackage{setspace}
\setstretch{1.3}

% Paragrafi
\setlength{\parindent}{0pt}
\setlength{\parskip}{5pt}

% Miglioramento giustificazione
\emergencystretch=2em

% Collegamenti
\usepackage[hidelinks]{hyperref}

% Titolo
\title{Ontologia, principi filosofici}
\author{Luca Mandrelli}
\date{Versione 0.1 \\ \today}

\begin{document}

\maketitle

\begin{abstract}
% Qui inseriremo la sintesi dell'articolo.
Il presente articolo si propone di esaminare i principi
filosofici fondamentali dell'ontologia, con particolare
attenzione ai concetti di essere, ente, identità e relazione.
L'obiettivo è individuare le categorie concettuali attraverso
cui è possibile descrivere ed interpretare la realtà.
\end{abstract}

\newpage
\tableofcontents
\newpage

\section{Introduzione} %Capitolo 1
% Inizia qui il capitolo dell'articolo.
La parola ontologia deriva dal greco \textit{onto} (= ente) e \textit{logos} (= studio) %corsivo
e come tale rappresenta lo studio di cosa esiste. Tale disciplina filosofica di predispone 
di rispondere alla domanda: \textbf{"che cosa c'è ?"}. Differisce pertanto dalla metafisica, 
dal greco \textit{meta} (= oltre) e \textit{phisica} (= la fisica),
che si predispone pertanto di rispondere alla domanda: \textbf{"cosa è ?"}, cercando di spiegare
la ragione ultima, il perchè, il principio e il funzionamento di quello che esiste.
Può essere considerata come il capitolo preliminare alla metafisica.
\\ L'ontologia \textbf{formale} studia l'ente, o entità, in quanto tale e le strutture logiche astratte
che valgono per qualunque ente: oggetto, proprietà, identità, parte/tutto, dipendenza, 
relazione. 
\\ L'ontologia \textbf{materiale} studia gli aspetti specifici e il contenuto concreto di un ambito
della realtà.

\subsection{La priorità dell'ontologia e i suoi limiti} %Paragrafo 1.1
Data la sua natura filosofica nasce l'obiezione che debba essere
dell'ontologia il compito di definire se \textit{qualcosa} esista davvero.
D'altra parte può risultare complesso definire il confine tra il compito di definire
cosa esiste \textit{(rif.ontologia)} e cosa è \textit{(rif.metafisica)}. Ad esempio, nelle scienze
naturali è impossibile definire una nuova specie senza definire cosa essa è. 
\\ Questo relegherebbe la disciplina ontologica ad un mero compito di classificazione di ciò che è definito a priori.
\\ Tuttavia l'inventario non va inteso come un censimento di tutti i cittadini del mondo, quanto
nella scelta e definizione della lista dei tipi di entità mda includere in tale catalogo. 
\\ Tali categorie hanno una valenza generale riferendosi a tipi quali \textit{classe, proprietà, proposizione}.
\\ Se specificare le voci dell'inventario non debba fornire l'estensione completa, dovrebbe infatti 
favorirne almeno \textit{l'intensione}, ovvero quella definizione intrinseca che è basata sulle 
caratteristiche. 
\\Si pensi a tal fine alla matematica e quelle entità come il numero naturale: dal
punto di vista ontologico si possono intendere come quegli individui accomunati dall'insieme delle
proprietà che definiscono il corpus assiomatico dell'aritmetica, indipendentemente dalle proprietà
insiemistiche, che ne riflettono la loro natura più "metafisica".

\subsection{La priorità della metafisica e i suoi limiti} %Paragrafo 1.2
Ribaltando l'ordine delle priorità ci si potrebbe riferire all'assunto che prima si debba specificare 
che cosa possa esistere e quindi cosa esiste davvero. Questo assunto mette in priorità la questione
metafisica rispetto a quella ontologica.
\\ Tuttavia al fine di una ricerca ontologica priva di "buchi" bisognerebbe pensare in grande, fissando
l'intensione del maggior numero di categorie possibili e quindi rivolgersi alla loro estensione.
\\ Permane quindi il problema di verificare l'esaustività delle categorie proposte. Inoltre
esiste anche la questione delle entità miste, rappresentate dalle somme mereologiche che rende improbo
il compito di definire a priori in maniera completa le categorie metafisiche. 

\subsection{L'analogia tra ontologia e metafisica, lessico e logica} %Paragrafo 1.3
Volendo stabilire una complementarietà tra ontologia e metafisica nella risposta alla medesima domanda
\textbf{"che cosa esiste"} si può sfruttare la seguente analogia:
\begin{itemize}
    \item l'ontologia risponde alla domanda in virtù delle costanti non logiche, siano queste nomi, aggettivi, etc.
    \item la metafisica risponde alla domanda in virtù dell'esplicitazione di assunti esistenziali convogliati
    da un enunciato in relazione alla sua forma logica.
\end{itemize}

\subsection{L'ontologia attraverso altre interpretazioni} %Paragrafo 1.4
Secondo \textbf{\textit{Heidegger}} l'indagine ontologica rappresenta un'analisi dei modi di darsi dell'essere
ad un soggetto, riportando la stessa ad una visione fenomenologica della realtà.
\\ Secondo \textbf{\textit{Hartmann}} l'ontologia è stratificata in livelli differenti di realtà, 
\textit{fisico-materiale, organico, psichico e spirituale}. La sintesi attraverso un minimo comune denominatore
sarebbe impossibile e pertanto occorre ricondursi a strati differenti per descrivere l'essenza dell'essere. 

\subsection{Ontologia formale vs materiale} %Paragrafo 1.4
L'ontologia formale studia le strutture ultime in cui è organizzata la realtà. In tal senso si contrappone
a quella materiale che si occupa di specifici settori della realtà (es. fisica, biologia, etc\dots). Si occupa pertanto
dell'essere indipendentemente dalle sue manifestazioni contingenti. Non si predispone pertanto di redigere un inventario
dettagliato, nè di specificarne la natura, ma si occupa di definire ciò che esiste sotto il profilo delle sue 
caratteristiche necessarie in quanto definite in ogni dominio.
\\Vista come \textbf{algebra dell'essere} non si occupa di quali entità esistano o del perchè esistano, quanto dei
nessi che sussistono tra le entità in virtù delle loro condizioni di esistenza. Definisce pertanto le leggi
\begin{itemize}
    \item teoria dell'identità e transitività della relazione di identità
    \item relazione di dipendenza
    \item appartenenza mereologica (parte-tutto)
\end{itemize}
Il legame tra ontologia formale e logica porta all'idea che queste possano rappresentare due aspetti di una scienza
unificata che si rivolge a tutte le forme dell'essere, determinandone lo studio delle diverse modalità dell'essere
sotto leggi generali.

\section{Come si fa ontologia?} %Capitolo 2
% Inizia qui il capitolo dell'articolo.
Verranno di seguito esposti gli strumenti analitici di cui ci si può servire per fare ontologia.

\subsection{L'analisi del linguaggio e i suoi limiti} %Paragrafo 2.1
Per alcuni filosofi di orientamento analitico l'analisi del linguaggio rappresenta una delle poche strade per
raggiungere risultati concreti nell'ambito dell'ontologia formale.
\\L'ontologia formale consiste nella ricerca di quelle cose che rendono veri gli enunciati veri del linguaggio
o che esistono affinchè gli enunciati possano essere veri.

\subsection{Verità ed esistenza} %Paragrafo 2.2
Si definisce \textit{impegno ontologico} il riconoscimento dell'essere in quanto essere come valore di una
variabile vincolata. L'impegno ontologico è soddisfatto quando viene definito un giusto livello di analisi
semantica.
\begin{equation}
    \mbox{Esiste $x$ tale che $x$ è un oggetto.}
    \label{eq:esistenza_oggetto}
\end{equation}
È possibile rivelare l'impegno ontologico formulando gli enunciati nella forma della logica del primo ordine. Per chiarire 
questo concetto prendiamo due enunciati e analizziamo come differiscano tra loro in relazione all'impegno ontologico.
\begin{equation}
    \exists x \; (Corvo(x))
    \label{eq:impegno_ontologico}
\end{equation}
L'enunciato \eqref{eq:impegno_ontologico} rivela la condizione di almeno un x ad essere corvo. Questo enunciato risulta
profondamente diverso dal seguente:
\begin{equation}
    \forall x \; (Corvo(x) \rightarrow Nero(x))
    \label{eq:quant_universale}
\end{equation}
in cui il fatto che i corvi siano tutti neri non definisce l'impegno ontologico dell'esistenza dei corvi.
Come si fa a riconoscere il giusto livello di analisi semantica?
Secondo \textbf{\textit{Berti}} se è dubbio che la metafisica si riduca alla fisica, è altrettanto dubbio che 
l'ontologia si riduca alla linguistica. Prendiamo due enunciati di questa natura:
\begin{equation}
    \mbox{C'è un taglio in questa tela}
    \label{eq:taglio_tela}
\end{equation}
\begin{equation}
    \mbox{Questa tela è tagliata}
    \label{eq:tela_tagliata}
\end{equation}
L'enunciato \eqref{eq:taglio_tela} presuppone l'esistenza dell'entità "taglio" che per \eqref{eq:tela_tagliata}
rappresenta una allucinazione ontologica, nella misura in cui all'interno dell'enunciato stesso non si riconosce
l'entità "taglio" ma solo quella di "tela".
\\ Entrambi concordano sulla verità dell'affermazione ma non su cosa la renda vera. Così il filosofo \eqref{eq:taglio_tela}
può tacciare l'altro di miopia ontologica per non aver annoverato i tagli nell'inventario delle entità.
\\ Basandosi sulla sola \textit{linguistica} non è perciò chiaro quale sia la direzione di parafrasi che cattura 
la forma logica della verità:
\ E' la presenza del taglio che rende la tela tagliata oppure essendo la tela tagliata c'è un taglio?

\subsection{Significato ed inferenza} %Paragrafo 2.3
Il significato di un enunciato è determinato in buona parte anche dalle relazioni logiche che lo legano ad altri
enunciati. Torniamo a tal fine ad esaminare degli enunciati derivati da quelli del paragrafo precedente.
\begin{equation}
    \mbox{C'è un taglio diagonale in questa tela}
    \label{eq:taglio_diag_tela}
\end{equation}
\begin{equation}
    \mbox{Questa tela è tagliata in diagonale}
    \label{eq:tela_tagliata_diag}
\end{equation}

L'enunciato \eqref{eq:taglio_diag_tela} è proprio di un atteggiamento filosofico \textit{induzionista} a differenza
del \eqref{eq:tela_tagliata_diag} che è di stampo prettamente \textit{eliminativista}. Analizzando sul profilo logico
 l'enunciato \eqref{eq:taglio_diag_tela}, se esso è vero allora lo è implicitamente anche \eqref{eq:taglio_tela}. Viceversa 
non esistono leggi logiche che instaurano nessi implicativi tra \eqref{eq:tela_tagliata_diag} ed \eqref{eq:tela_tagliata}.
Tuttavia, di fatto, tale inferenza non è semantica ma logica.
\\ Questo evidenzia un limite nella definizione di cosa la \textit{logica classica} contempli e cosa no (la comune 
\textit{logica dei predicati} non si occupa in maniera esplicita degli avverbi).

\subsection{Rivelare e stipulare} %Paragrafo 2.4
Gli approcci finora tenuti conducono ad impasse riconducibili alla forma linguistica e al linguaggio. 
Per uscirne viene in ausilio l'utilizzo di parafrasi a favore del definire l'impegno ontologico di una asserzione.
Si distinguono due approcci nella esplicitazione del'impegno ontologico, rispettivamente \textit{ermeneutico-interpretativo}
e \textit{stipulativo}
\\ Nella forma \textit{ermeneutica} la parafrasi è tesa a rivelare la forma logica e le effettive condizioni di verità.
\\ Nella forma \textit{stipulativa} la parafrasi è usata allo scopo di evitare fraintendimenti in quanto si considera
il linguaggio ordinario \textit{ontologicamente neutrale}, pertanto la parafrasi rappresenta esclusivamente l'interpretazione
che si tende adottare e che risulti equivalente all'enunciato di partenza.

\subsection{Prescrivere e descrivere} %Paragrafo 2.5
L'attributo prescrittivo all'ontologia deriva da una posizione in merito alla quale si vorrebbe che l'altro accetti
le proprie teorie o quindi la propria visione del mondo.
\\ Dall'altra parte la concezione descrittiva per contro si riferisce ad una analisi ontologica che descrive l'ossatura
del nostro pensiero sul mondo, ovvero la rappresentazione dello stesso nel nostro sistema cognitivo \textbf{\textit{Kant}}.
Il valore della concezione descrittiva dell'ontologia risiede quindi nell'acquisire la padronanza dei concetti che usiamo
per descrivere il mondo.
\\ Sul lato pratico, secondo \textbf{\textit{Hayes}} e la sua \textit{Fisica Ingegnua} un valore oggettivo della concezione
descrittiva può essere quello di dotare un robot delle categorie atte a filtrare efficaciemente la realtà in cui si deve
operare.
\\ L'\textit{ontologo prescrittivo} quindi considera il linguaggio naturale come qualcosa con cui occorre fare i conti
nel corso dell'indagine ontologica. Alla domanda "che cosa esiste" risponde guardando quali siano le entità che debbano
esistere per far sì che una serie di fatti e di intuizioni siano rispettate. Il rendere conto alle condizioni di verità
degli enunciati del linguaggio viene fatto dopo la messa a fuoco delle convinzioni ontologiche. Questo è quanto va inscritto
nel concetto di \textit{analisi stipulativa}.
\\ L'\textit{ontologo descrittivo}, dall'altra parte, è colui per il quale il linguaggio naturale si erige a principale
fonte di ispirazione. E' tramite il linguaggio che il sistema dei concetti attraverso i quali pensiamo il mondo trova la sua
espressione più esplicita. Sulla base di questo i fondamenti e il campo di azione dell'indagine ontologica finiscono con 
il concidere con la teoria del significato.
\\ Per questo l'\textit{ontologo descrittivo} finirà per riconoscersi con asserzioni del tipo \eqref{eq:taglio_tela} rifiutando
l'\textit{eliminativismo} più spinto, poichè le entità che l'asserzione definisce sono rappresentate nel lessico del nostro
linguaggio.

\subsection{Assoluto e relativo} %Paragrafo 2.6
L'accezione \textit{assoluta} mira ad una determinazione universale ed univoca dell'inventario del mondo e per questo 
è legata ad approcci di tipo \textit{prescrittivo}. Tuttavia il legame tra \textit{ontologia prescrittiva} e \textit{realismo}
non è così forte. Il \textit{realismo} non esclude che ogni teoria debba formularsi all'interno di un opportuno sistema di
coordinate, quanto piuttosto il fatto che esistano dei sensi vietati di rappresentazione i quali obbligano gli schemi del mondo 
a rispettare alcune teorie ammissibili.
\\ L'accezione \textit{relativa}, dall'altra parte, punta ad ricorrere ad un opportuno sistema di coordinate. Da questo 
punto di vista si osserva come l'approccio descrittivo all'ontologia sia naturalmente legato ad elaborazioni in chiave 
\textit{relativistica}.

\end{document}
